\documentclass[12pt]{article}
\usepackage[utf8]{inputenc}
\usepackage[T1]{fontenc}
\usepackage{hyperref}
\title{Sara Kaba Language Resources for K-12 Curriculum Development}
\author{Compiled by Claude Sonnet 5.5 (claude-sonnet-5-5) for Andr\'as Kornai}
\date{}
\begin{document}
\maketitle

\section{Introduction}
Sara Kaba (ISO 639-3 sbz) is a Central Sudanic (Bongo--Bagirmi) language of the Central African Republic. The names ``Kaba'' and ``Sara Kaba'' are used loosely. The Kaba languages are three to five languages in Chad and the Central African Republic: Kaba proper (Ta Sara), Kaba Na, Kaba Deme, Kaba Dunjo and Kulfa. They should not be confused with the Sara languages proper, and sources disagree on which varieties are distinct languages. Wikitongues lists sbz as unwritten and vigorous (EGIDS 6a). The searches found no school materials. This dossier is short.

\section{Primary Educational Resources}
No primers, textbooks or school programmes for sbz were found.

\section{Linguistic and Reference Materials}
\begin{itemize}
\item Wikipedia, Kaba language (sbz): \url{https://wikipedia.com/wiki/ISO_639:sbz}
\item Wikipedia, Kaba languages (overview of the group): \url{https://en.wikipedia.org/wiki/Kaba_languages}
\item Wikipedia, Kabba language (ksp), which cites Rosmarie Moser, \emph{Kabba: A Nilo-Saharan Language of the Central African Republic} (Lincom, 2004), as the main description of Kabba: \url{https://en.wikipedia.org/wiki/Kabba_language}
\item BnF catalogue heading for Kaba (langue), which lists related headings for Kaba dictionaries and tales; I could not see the titles: \url{https://data.bnf.fr/12050495/kaba__langue_/fr.pdf}
\end{itemize}
Caution: Moser's book and the ksp entry may describe a different Kaba variety than sbz. I did not verify this. The Dictionary of Na (Keegan and Koutou, 2015), mentioned in a search summary, covers the related Na variety, and I did not see a URL for it.

\section{Religious and Community Resources}
None found for sbz specifically.

\section{Digital and Audio Resources}
Wikitongues language page for sbz: \url{https://wikitongues.org/languages/sbz/}. Wikitongues also has a page for Kaba (ksp): \url{https://wikitongues.org/languages/ksp/}.

\section{Curriculum Development Context}
Very little is documented for sbz. Work would begin with deciding which Kaba variety is meant, then documenting it and developing an orthography with the community. The Kaba literature from related varieties may help, but its fit to sbz should be checked.
\end{document}